\documentclass[10pt,a4paper]{amsart}
\usepackage[margin=19mm]{geometry}
\usepackage{amsmath,amssymb,mathtools}
\usepackage[scr=rsfs,cal=euler]{mathalfa}
\usepackage{xfrac}
\usepackage[unicode,hidelinks,bookmarks=false]{hyperref}
\newcommand{\ie}{i.e.\ }
\newcommand{\cf}{cf.\ }
\newcommand{\resp}{respectively\ }
\newcommand{\Subsection}{\subsection}
\providecommand{\nobreakdash}{\nobreak\hbox{-}}
\setlength{\emergencystretch}{2em}
\multlinegap0pt
\usepackage{bm}
\usepackage{bbm}
\usepackage{dsfont}
\usepackage{xspace}
\usepackage{cancel}
\usepackage{mathtools}
\usepackage{amsthm}
\makeatletter
\@ifundefined{proposition}{\newtheorem{proposition}{Proposition}}{}
\makeatother
\title[When Is a Relevance Threshold Statistically Resolvable? Mini]{When Is a Relevance Threshold Statistically Resolvable? Minimax Limits for Effect Classification}
\author{Subir Hait}
\date{Source: arXiv:2609.15811v2 (CC BY 4.0)}
\begin{document}
\maketitle

\noindent\emph{Source and licence.} When Is a Relevance Threshold Statistically Resolvable? Minimax Limits for Effect Classification. Version arXiv:2609.15811v2 (CC BY 4.0), distributed under the Creative Commons Attribution 4.0 International licence, as recorded in the arXiv licence metadata for this exact version and shown on the abstract page https://arxiv.org/abs/2609.15811v2. Full rights notice: licenses/arxiv-2609.15811-CC-BY-4.0.txt.\par\smallskip

\noindent\textit{Editorial scope.} Counted unit: the Proposition whose source label is \texttt{prop:critical-exact} in the article (source0.tex lines 393--421, source label \texttt{prop:critical-exact}), delivered together with the article's own complete original proof of it (source0.tex lines 423--450, one \texttt{proof} environment of 28 lines). No mathematics is written, repaired or re-derived: statement and proof are the article's own text line for line. Required context retained uncounted: none. Every definition, display and label of those lines is retained unchanged. Omitted from this package and not redistributed: 1--392, 422--422, 451--1152. Nothing inside a retained excerpt is omitted or altered: every retained body is delivered exactly as the article's lines are written, so no omission receipt of any kind accompanies this package. Citations: the retained excerpts carry no citation command at all, so no bibliography entry of the article is transcribed and no reference is redistributed. Reference adaptation: none was needed and none was made; every reference inside the delivered unit resolves inside it, so no unresolved reference or citation is produced. Printed slips kept literal: no printed word, symbol, hypothesis or number of the counted statement or of its proof was altered. Layout and class: the article's own class is \texttt{article} with packages including geometry, amsmath, amsthm, mathtools, newtxtext, newtxmath, graphicx, booktabs, microtype, setspace, biblatex, hyperref, enumitem, caption; the wrapper is the lane's pinned \texttt{amsart} editorial wrapper at 10pt on A4, which restates the theorem-like environments the retained text needs and adds the article's own macro definitions that the retained text needs (none), with extra wrapper packages (bm, bbm, dsfont, xspace, cancel, mathtools). The delivered numbering is the unit's own. Assets: no retained span contains a figure environment, a graphics-inclusion command, a PDF-page inclusion, a table or a TikZ picture; no image file is copied into this package, the package declares no asset dependency and no image file is redistributed with it. Licence: version arXiv:2609.15811v2 of this article is distributed under the Creative Commons Attribution 4.0 International licence, as recorded in the arXiv licence metadata for this exact version and shown on the abstract page https://arxiv.org/abs/2609.15811v2. A full rights notice is delivered at licenses/arxiv-2609.15811-CC-BY-4.0.txt. Characters: every retained excerpt is pure ASCII, so the delivered engine, XeLaTeX with its default Latin Modern text font, reports no missing character.
\begin{proposition}[Exact minimax characterization in the critical Gaussian experiment]
\label{prop:critical-exact}
Let $g_m$ denote the equal mixture of $N(m,1)$ and $N(-m,1)$,
\[
g_m(y)=\frac12\{\varphi(y-m)+\varphi(y+m)\}
=\varphi(y)e^{-m^2/2}\cosh(my).
\]
There is a unique $c^*=c^*(\kappa,\varepsilon)>0$ satisfying
\begin{equation}
e^{-b^2/2}\cosh(bc^*)
=
e^{-a^2/2}\cosh(ac^*).
\label{eq:cstar}
\end{equation}
The rule
\begin{equation}
\phi^*(Y)=\mathbf1\{|Y|>c^*\}
\label{eq:phistar}
\end{equation}
is minimax for $\mathcal N_\infty$ versus $\mathcal M_\infty$. Its exact risk is
\begin{equation}
\begin{aligned}
R_\infty^*(\kappa,\varepsilon)
={}&\Phi(a-c^*)+\Phi(-a-c^*)\\
&+\Phi(c^*-b)-\Phi(-c^*-b).
\end{aligned}
\label{eq:exactrisk}
\end{equation}
\end{proposition}

\begin{proof}
For any rule $\phi$, the worst-case risk dominates the average risk under equal boundary mixtures:
\[
R_\infty(\phi)
\ge
E_{g_a}\phi+E_{g_b}(1-\phi).
\]
The infimum of the right-hand side is the simple-versus-simple Bayes risk for $g_a$ against $g_b$, attained by the likelihood-ratio rule that chooses the meaningful class when $g_b(y)>g_a(y)$.

For $x=|y|$,
\[
\frac{g_b(y)}{g_a(y)}
=
e^{-(b^2-a^2)/2}
\frac{\cosh(bx)}{\cosh(ax)}.
\]
For $x>0$, the derivative of its log is
\[
b\tanh(bx)-a\tanh(ax)>0,
\]
because $m\mapsto m\tanh(mx)$ is strictly increasing for $m>0$. At $x=0$ the ratio is smaller than one, whereas it diverges as $x\to\infty$. Hence there is a unique crossing $c^*>0$, characterized by \eqref{eq:cstar}, and the Bayes rule is \eqref{eq:phistar}.

It remains to verify that the boundary mixtures are least favorable for the composite problem. For any fixed $c>0$ and $\mu\ge0$,
\[
P_\mu(|Y|>c)=\Phi(\mu-c)+\Phi(-\mu-c)
\]
is increasing in $\mu$. Therefore the largest false-positive probability over $|\mu|\le a$ occurs at $|\mu|=a$, and the largest false-negative probability over $|\mu|\ge b$ occurs at $|\mu|=b$. For the likelihood-ratio threshold $c^*$, the composite worst-case risk consequently equals the boundary-mixture Bayes risk. The lower bound above is therefore attained, proving minimaxity. Substituting the two boundary error probabilities yields \eqref{eq:exactrisk}.
\end{proof}

\end{document}
